\documentclass[10pt]{article}

\usepackage[utf8]{inputenc}

\usepackage[T1]{fontenc}

\usepackage{amsmath}

\usepackage{amsfonts}

\usepackage{amssymb}

\usepackage[version=4]{mhchem}

\usepackage{stmaryrd}

\usepackage{bbold}



\begin{document}

HAKЛОН ди фр акцион ного кону са - см. Редже полюсов метод.



НАКЛОННОЙ ПРОИЗВОДНОЙ ЗАДАЧА - сМ. КРаевая задача теории потенциала.



HAKPЫTИE - расслоение с дискретным слоем, или локально топологическое отображение $f: X \rightarrow Y$ (у каждой точки $y \in Y$ существует такая окрестность, что $f^{-1}(U)$ состоит из открытых множеств, гомеоморфных $U$ ). $\quad$. И. Войцеховский. НАПРАВЛЕНИЙ КРИВИЗНА - сМ. Вейля кривизна.



НАПРАВЛЕНИЙ ПОЛЕ - область На плоскости, в Каждой точке которой задана прямая. Для задания таких прямых достаточно знать тангенс угла между прямой и осью $O x$, поэтому понятие Н. п. чаще всего связывается с Н. п. касательных нек-рого семейства кривых, заполняющих данную область. В частности, через каждую точку $(t, x)$ области определения $G \subset \mathbb{R}^{2}$ правой части уравнения



\[

\frac{d x}{d t}=f(t, x)

\]



можно провести отрезки малой длины с угловым коэффициентом $f(t, x)$, так что (ориентированный) угол между осью $O t$ и этим отрезком равен $\operatorname{arctg} f(t, x)$. Часто дифференциальное уравнение (1) рассматривается в совокупности с дифференциальным уравнением



\[

\frac{d t}{d x}=F(t, x),

\]



где $F(t, x)=1 / f(t, x)$ для точек $(t, x) \in G$, в к-рых $f(t, x) \neq 0$, и $F(t, x)=0$ для точек $(t, x) \in \mathbb{R}^{2} \backslash G$, в к-рых функцию $F(t, x)$ можно доопределить этим значением по непрерывности. Тем самым для пары уравнений $(1),(2)$ область $G$ расширяется до области $G_{0}$ за счет пополнения точками, в к-рых направление параллельно оси $O x$, а интегральным кривым разрешается иметь и точки с вертикальной касательной. Если в области $G$ для уравнения (1) (или в области $G_{0}$ для пары уравнений (1), (2)) изобразить Н. п. достаточно подробно, то по построенным отрезкам можно составить примерное качественное представление о картине поведения интегральных кривых. Это соображение лежит в основе приближенного графич. метода решения уравнения (1) метода изоклин, в к-ром построение Н. п. осуществляется с помощью изоклин.



В. И. Битюиков.



НАПРЯЖЕНИЙ ТЕНЗОР - тензор, определяющиЙ распределение внутренних сил деформированного твердого тела. В силу закона сохранения импульса:



\[

\frac{\partial \Pi_{i}}{\partial t}+\frac{\partial \sigma_{i k}}{\partial x_{k}}=0, \quad i, k=1,2,3,

\]



где П - плотность импульса, Н. т. $\sigma_{i k}$ может быть определен как поток $i$-й составляющей плотности импульса через единичную площадку, перпендикулярную $k$-й оси.



Если $\boldsymbol{F}$ - суммарная сила, действующая на единицу объема тела, то Н. т. определяется соотношением, являющимся следствием (1):



\[

\int F_{i} d v=\int \sigma_{i k} d f_{k},

\]



где $d v$ - элемент объема, $d f$ - элемент площади.



Наиболее простой вид Н. т. имеет в случае равномерного всестороннего сжатия тела давлением $p$, а именно:



\[

\sigma_{i k}=-p \delta_{i k}

\]



где $\delta_{i k}$ - символ Кронекера.



И. В. Барьяхтар. НАТУРАЛЬНЫЙ ПАРАМЕТР н а к р и в о Й $x(t)$ в рИмановом многообразии $\left(M^{n}, g_{i j}\right)$ - длина $s(t)$ кривой $x(t)$, отсчитываемая от фиксированной точки $p_{0}=x(a)$ до текущей точки $p=x(t)$, то есть в пределах одной карты



\[

s(t)=\int_{a}^{t} \sqrt{g_{i j} \frac{\partial x^{i}}{\partial t} \frac{\partial x^{j}}{\partial t}} d t

\]



где $\left(x^{1}, \ldots, x^{n}\right)$ - локальная система координат в рассматриваемой карте многообразия $M^{n}$. Любые два натуральных параметра $s, s^{\prime}$ отличаются друг от друга на нек-рое ортогональное преобразование прямой, то есть $s^{\prime}= \pm s+$ const, так как можно изменить направление вдоль кривой на противоположное и сдвинуть начальную точку отсчета.



Лит.: [1] Дубровин Б.А., Новиков С.П., Фоменко А.Т., Современная геометрия. Методы и приложения, М., 1986.



В. В. Трофимов.



НАЧАЛЬНАЯ ФАЗА к о Л е б а н и Я - см. Гармоника.



НАЧАЛЬНОЕ ВОЗМУЩЕНИЕ - СМ. Устойчивость движения.



НАЧАЛЬНЫЕ УСЛОВИЯ - см. Математической физики классические уравнения, Краевая задача.



НАЧАЛЬНЫЙ МОМЕНТ - сМ. КвантовЫй случайнЫй проuecc.



НЕАДИАБАТИЧЕСКИЙ ПЕРЕХОД - см. Адиабатический инвариант.



НЕБЛУХДАЮЩАЯ ТОЧКА - точКа для отображения $f$, для любой окрестности $U$ которой найдется такое $n$, что



\[

f^{n}(U) \cap U \neq \varnothing .

\]



Точка, не обладающая этим свойством, называется б луж д а ю щ е й. Множество Н. т. замкнуто и инвариантно относительно $f$. Аналогичное понятие можно ввести и для динамич. систем с непрерывным временем, то есть потоков.



Я. Б. Песин.



НЕВАНЛИННЫ КЛАСС $N$ в единичном круге $|\zeta|<1$ - класс, состоящий из всех функций вида:



\[

f(\zeta)=\Phi(\zeta) / \Psi(\zeta),

\]



где $\Phi(\zeta), \Psi(\zeta)$ - ограниченные аналитические в круге $|\zeta|<1$ функции. Те орема Неванлинны: функция $f(\zeta)$, аналитическая в круге $|\zeta|<1$, принадлежит Н. к. тогда и только тогда, когда



\[

\sup \left\{\int_{0}^{2 \pi} \ln ^{+}\left|f\left(r e^{i \theta}\right)\right| d \theta ; \quad 0 \leqslant r<1\right\}<\infty .

\]



Н. к. в верхней полуплоскости $z=x+i y, y>0$ определяется преобразованием Кэли



\[

\zeta \mapsto z=i(1+\zeta) /(1-\zeta) .

\]



Класс $N$ введен Р. Неванлинной (R. Nevanlinna) в конце 20-х гт. 20 в. $\quad$ B. В. Жаринов. НЕВЛОХИМОСТИ TEOPEМА - см. Симплектическая геометрия.



НЕВОЗВРАТНОЕ СЛУЧАЙНОЕ БЛУЖДАНИЕ см. Случайное блуждание.



НЕВОЗМУЩЕННОЕ ДВИХЕНИЕ - см. Устойчивость движения.



НЕВЫРОЖДЕННАЯ ТОЧКА ПЕРЕВАЛА - см. Перевала метод.



НЕВЫРОЖДЕННОСТИ УСЛОВИЕ Гам ил ьтон о в о Й с и с т е м ы-см. Интегрируемая гамильтонова система; невырожденность.



НЕВЫРОЖДЕННЫЙ ЛАГРАНЖИАН - СМ. ВариационНое исчисление в целом.



НЕГОЛОНОМНАЯ СвяЗь, д и фф е рен ци аль ная с в я зь, кине мат и ч е ская с в я зь,- связь неголономной системы. См. также Связи механические.



НЕГОЛОНОМНАЯ СИСТЕМА - механическая система, на которую кроме геометрических наложены еще дифференциальные (кинематические) связи. Связи системы называются дифференциальными (неголономным и), если они накладывают ограничения не только на координаты, но и на скорости точек и выражаются математически уравнениями вида:



\[

f_{k}\left(x_{i}, y_{i}, z_{i}, \dot{x}_{i}, \dot{y}_{i}, \dot{z}_{i}, t\right)=0, \quad k=1,2, \ldots, r,

\]



где $x_{i}, y_{i}, z_{i}-$ координаты, $\dot{x}_{i}, \dot{y}_{i}, \dot{z}_{i}-$ проекшия скоростей, $t$ - время, $r$ - число наложенных связей.



НЕГОЛОНОМНАЯ 375





\end{document}